\setcounter{section}{5}






  \zwischenueberschrift{Übungsaufgaben}

\inputaufgabe
{}
{





Es sei

\mavergleichskette
{\vergleichskette
{ F
}
{ \in }{ K[X_1 , \ldots ,  X_n]
}
{  }{
}
{    }{
}
{   }{
}
}
{}{}{}
ein
\definitionsverweis {homogenes Polynom}{}{}
mit
\definitionsverweis {Nullstellenmenge}{}{}
$V(F)$. Zeige, dass für jeden Punkt

\mavergleichskette
{\vergleichskette
{ P
}
{ \in }{ V(F)
}
{  }{
}
{    }{
}
{   }{
}
}
{}{}{}
und jeden Skalar

\mavergleichskette
{\vergleichskette
{ \lambda
}
{ \in }{ K
}
{  }{
}
{    }{
}
{   }{
}
}
{}{}{}
auch

\mavergleichskette
{\vergleichskette
{ \lambda P
}
{ \in }{ V(F)
}
{  }{
}
{    }{
}
{   }{
}
}
{}{}{}
ist.





}
{} {}

\inputaufgabe
{}
{





Bestimme die Faktorzerlegung für die Polynome

\mavergleichskettedisp
{\vergleichskette
{ X^n-Y^n
}
{ \in} { {\mathbb C}[X,Y]
}
{ } {
}
{ } {
}
{ } {
}
}
{}{}{}
für

\mavergleichskette
{\vergleichskette
{ n
}
{ \in }{ \N_+
}
{  }{
}
{    }{
}
{   }{
}
}
{}{}{.}





}
{} {}

\inputaufgabegibtloesung
{}
{





Es sei $K$ ein
\definitionsverweis {algebraisch abgeschlossener Körper}{}{}
und sei

\mavergleichskette
{\vergleichskette
{F
}
{ \in }{ K[X,Y]
}
{  }{
}
{    }{
}
{   }{
}
}
{}{}{}
ein
\definitionsverweis {homogenes Polynom}{}{.}
Zeige: $F$ zerfällt in
\definitionsverweis {Linearfaktoren}{}{.}





}
{} {}

\inputaufgabe
{}
{





Es sei $K$ ein
\definitionsverweis {Körper}{}{}
und

\mavergleichskette
{\vergleichskette
{ F
}
{ \in }{ K[X_1 , \ldots ,  X_n]
}
{  }{
}
{    }{
}
{   }{
}
}
{}{}{}
ein
\definitionsverweis {homogenes Polynom}{}{.}
Es sei

\mavergleichskette
{\vergleichskette
{ F
}
{ = }{ GH
}
{  }{
}
{    }{
}
{   }{
}
}
{}{}{}
eine Faktorzerlegung. Zeige, dass
\mathkor {} {G} {und} {H} {}
ebenfalls homogen sind.





}
{} {}

\inputaufgabe
{}
{





Es sei $R$ ein
\definitionsverweis {kommutativer Ring}{}{}
und

\mavergleichskette
{\vergleichskette
{ P
}
{ = }{ R[X_1 , \ldots , X_m]
}
{  }{
}
{    }{
}
{   }{
}
}
{}{}{}
der
\definitionsverweis {Polynomring}{}{}
darüber in $m$ Variablen. Es sei

\mavergleichskette
{\vergleichskette
{ {\mathfrak m}
}
{ = }{ \left(  X_1 , \ldots , X_m  \right)
}
{  }{
}
{    }{
}
{   }{
}
}
{}{}{}
das von den Variablen
\definitionsverweis {erzeugte Ideal}{}{.}
Zeige, dass

\mavergleichskette
{\vergleichskette
{ {\mathfrak m}^n
}
{ = }{ P_{\geq n}
}
{  }{
}
{    }{
}
{   }{
}
}
{}{}{}
ist, wobei $P_{\geq n}$ das Ideal in $P$ bezeichnet, das von allen
\definitionsverweis {homogenen Polynomen}{}{}
vom Grad $\geq n$ erzeugt wird.





}
{} {}

\inputaufgabe
{}
{





Zeige, dass ein
\definitionsverweis {homogenes Polynom
}{}{}
unter einer linearen Variablentransformation homogen vom gleichen Grad bleibt, und dass dies bei einer
\definitionsverweis {affin-linearen Variablentransformation
}{}{}
nicht sein muss.





}
{} {}

\inputaufgabe
{}
{





Zeige, dass für
\definitionsverweis {affin-algebraische Mengen}{}{}

\mavergleichskette
{\vergleichskette
{ V,V'
}
{ \subseteq }{ { {\mathbb A}_{  K }^{  n   } }
}
{  }{
}
{    }{
}
{   }{
}
}
{}{}{}
die Beziehung der
\definitionsverweis {affin-linearen Äquivalenz}{}{}
eine
\definitionsverweis {Äquivalenzrelation}{}{}
ist.





}
{} {}

\inputaufgabe
{}
{





Es sei

\mavergleichskette
{\vergleichskette
{ P
}
{ = }{ (a,b)
}
{  }{
}
{    }{
}
{   }{
}
}
{}{}{}
ein Punkt in der affinen Ebene und $L$ und $L'$ verschiedene Geraden durch $P$. Es sei

\mathbed {C = V(F)} {}
 {F \in K[X,Y]} {}
 {} {} {} {,}
eine ebene algebraische Kurve. Beschreibe explizit eine Variablentransformation
\zusatzklammer {einen Koordinatenwechsel} {} {}
derart, dass in den neuen Koordinaten $P$ der Nullpunkt wird und die Geraden zum Achsenkreuz werden. Wie lautet die Kurvengleichung in den neuen Koordinaten?





}
{} {}

\inputaufgabe
{}
{





Es sei sowohl $C$ als auch $D$ eine ebene affin-algebraische Kurve, die jeweils aus der Vereinigung von drei
\zusatzklammer {verschiedenen} {} {}
Geraden bestehen, die sich jeweils in einem Punkt treffen. Zeige, dass es einen
\definitionsverweis {affin-linearen Koordinatenwechsel}{}{}
gibt, der $C$ in $D$ überführt.





}
{} {}

\inputaufgabe
{}
{





Es sei sowohl
\mathl{C}{} als auch $D$ eine ebene affin-algebraische Kurve, die jeweils aus der Vereinigung von vier
\zusatzklammer {verschiedenen} {} {}
Geraden bestehen, die sich jeweils in einem Punkt treffen. Zeige, dass es im Allgemeinen keinen
\definitionsverweis {affin-linearen Koordinatenwechsel}{}{}
gibt, der $C$ in $D$ überführt.





}
{} {}

Die folgenden beiden Aufgaben dienen dem Verständnis von

Satz 5.4

und

Korollar 5.5
.

\inputaufgabe
{}
{





Wende den Beweis zu

Satz 5.4

auf das Polynom $Y$
an.





}
{} {}

\inputaufgabe
{}
{





Wende
\zusatzklammer {den Beweis zu} {} {}

Satz 5.4

auf die Hyperbel
\mathl{XY-1}{} an.





}
{} {}

\inputaufgabe
{}
{





Wende
\zusatzklammer {den Beweis zu} {} {}

Satz 5.4

auf das Polynom

\mavergleichskette
{\vergleichskette
{ X^2Y^3+5X^3Y^2-X^2Y^2+3Y+7
}
{ \in }{  {\mathbb C}[X, Y]
}
{  }{
}
{    }{
}
{   }{
}
}
{}{}{}
an.





}
{} {}

\inputaufgabe
{}
{





Es sei

\mavergleichskette
{\vergleichskette
{ F
}
{ \in }{ {\mathbb C} [X,Y]
}
{  }{
}
{    }{
}
{   }{
}
}
{}{}{}
ein nicht-konstantes Polynom. Zeige, dass die zugehörige algebraische Kurve

\mavergleichskette
{\vergleichskette
{ C
}
{ = }{ V(F)
}
{  }{
}
{    }{
}
{   }{
}
}
{}{}{}
\definitionsverweis {überabzählbar}{}{}
viele Elemente besitzt.





}
{} {}

\inputaufgabegibtloesung
{}
{





Es sei

\mavergleichskette
{\vergleichskette
{M
}
{ = }{ \{P_1 , \ldots , P_n \}
}
{ \subseteq }{ K^2
}
{    }{
}
{   }{
}
}
{}{}{}
eine endliche Punktmenge in der Ebene über einem unendlichen
\definitionsverweis {Körper}{}{.}
\aufzaehlungzweiabc{Zeige, dass man $M$ als Durchschnitt von zwei
\definitionsverweis {algebraischen Kurven}{}{}
erhalten kann.
}{Zeige, dass man $M$ als Durchschnitt von zwei
\definitionsverweis {irreduziblen}{}{}
algebraischen Kurven erhalten kann.
}





}
{} {}

\inputaufgabe
{}
{





Berechne das Bild $\tilde{F}$
des Polynoms

\mavergleichskettedisp
{\vergleichskette
{F
}
{ =} { X^2Y+3XY-Y^3
}
{ } {
}
{ } {
}
{ } {
}
}
{}{}{}
unter dem durch

\mathdisp {X \longmapsto T^2+S-3,\, Y \longmapsto 3TS+S^2-T} {  }

definierten Einsetzungshomomorphismus
\maabbdisp {} {K[X,Y]  } { K[S,T]
} {.}





}
{} {}

\inputaufgabe
{}
{





Es sei $K$ ein unendlicher
\definitionsverweis {Körper}{}{}
und es sei

\mavergleichskette
{\vergleichskette
{ F
}
{ \in }{ K[X_1 , \ldots , X_n]
}
{  }{
}
{    }{
}
{   }{
}
}
{}{}{}
ein Polynom mit der zugehörigen Abbildung
\maabbdisp {F} { { {\mathbb A}_{  K }^{  n   } } } { { {\mathbb A}_{ K }^{  1   } }
} {.}
Zeige mit und ohne

Satz 5.10
,
dass das Bild von $F$ einpunktig oder unendlich ist.





}
{} {}

\inputaufgabe
{}
{





Es sei $K$ ein
\definitionsverweis {endlicher Körper}{}{}
mit $q$ Elementen und

\mavergleichskettedisp
{\vergleichskette
{ P_1 , \ldots , P_n
}
{ \in} { {\mathbb A}^{2}_{ K }
}
{ } {
}
{ } {
}
{ } {
}
}
{}{}{}
$n$ Punkte in der affinen Ebene. Zeige, dass es genau dann eine
\definitionsverweis {polynomiale Abbildung}{}{}
\maabbdisp {\varphi} { {\mathbb A}^{1}_{ K } } { {\mathbb A}^{2}_{ K }
} {}
mit

\mavergleichskette
{\vergleichskette
{ \operatorname{bild}  \varphi
}
{ = }{ \{ P_1 , \ldots , P_n \}
}
{  }{
}
{    }{
}
{   }{
}
}
{}{}{}
gibt, wenn

\mavergleichskette
{\vergleichskette
{ 1
}
{ \leq }{ n
}
{ \leq }{ q
}
{    }{
}
{   }{
}
}
{}{}{}
ist.





}
{} {}

\inputaufgabegibtloesung
{}
{





Man gebe ein Beispiel für eine polynomiale Abbildung
\maabbdisp {} { {\mathbb A}^{2}_{K} } { {\mathbb A}^{1}_{K}
} {}
derart, dass das Urbild von einem Punkt reduzibel ist, das Urbild von allen anderen Punkten aber irreduzibel.





}
{} {}

\inputaufgabegibtloesung
{}
{





Es sei $K$ ein
\definitionsverweis {Körper}{}{.}
Wir betrachten zu jedem

\mavergleichskette
{\vergleichskette
{ n
}
{ \in }{ \N_+
}
{  }{
}
{    }{
}
{   }{
}
}
{}{}{}
die Abbildung
\maabbeledisp {\varphi} {K^n } { K^n
} { \left(  \lambda_1   , \,   \lambda_2  , \,   \ldots , \,   \lambda_n                \right) } { \left( c_0   , \,  c_1  , \,   \ldots , \,  c_{n-1}                \right)
} {,}
die einem Nullstellentupel
\mathl{\left(  \lambda_1   , \,   \lambda_2  , \,   \ldots , \,   \lambda_n                \right)}{} das Koeffiziententupel
\mathl{\left( c_0   , \,  c_1  , \,   \ldots , \,  c_{n-1}                \right)}{}
\zusatzklammer {ohne die $1$} {} {}
des
\definitionsverweis {normierten Polynoms}{}{}

\mavergleichskettedisp
{\vergleichskette
{ (X- \lambda_1) (X- \lambda_2) \cdots (X- \lambda_n)
}
{ =} { P
}
{ =} { c_0 +c_1X + \cdots + c_{n-1}X^{n-1} +X^n
}
{ } {
}
{ } {
}
}
{}{}{}
zuordnet.
\aufzaehlungsieben{Beschreibe $\varphi$ explizit für

\mavergleichskette
{\vergleichskette
{n
}
{ = }{ 2
}
{  }{
}
{    }{
}
{   }{
}
}
{}{}{.}
}{Beschreibe $\varphi$ explizit für

\mavergleichskette
{\vergleichskette
{n
}
{ = }{ 3
}
{  }{
}
{    }{
}
{   }{
}
}
{}{}{.}
}{Begründe, dass die $\varphi$
\definitionsverweis {polynomiale Abbildungen}{}{}
sind.
}{Zeige, dass die
\definitionsverweis {Fasern}{}{}
von $\varphi$ endlich sind.
}{Wann ist die Faser zu einem Tupel
\mathl{\left( c_0   , \,  c_1  , \,   \ldots , \,  c_{n-1}                \right)}{} leer?
}{Was ist die maximale Anzahl in einer Faser? Man gebe Beispiele, dass diese Maximalanzahl für

\mavergleichskette
{\vergleichskette
{K
}
{ = }{\R
}
{  }{
}
{    }{
}
{   }{
}
}
{}{}{}
erreicht wird.
}{Es sei $K$ nun
\definitionsverweis {algebraisch abgeschlossen}{}{.}
Zeige, dass $\varphi$ surjektiv ist.
}





}
{} {}

\inputaufgabe
{}
{





Es sei
\maabbdisp {\varphi} { { {\mathbb A}_{  K }^{  r   } } } { { {\mathbb A}_{  K }^{  n   } }
} {}
eine
\definitionsverweis {polynomiale Abbildung}{}{}
und sei

\mavergleichskette
{\vergleichskette
{ T
}
{ \subseteq }{ { {\mathbb A}_{  K }^{  r   } }
}
{  }{
}
{    }{
}
{   }{}
}
{}{}{}
eine Teilmenge. Zeige, dass die Gleichheit

\mavergleichskettedisp
{\vergleichskette
{ \overline{\varphi(T)}
}
{ =} { \overline{\varphi(\overline{T} ) }
}
{ } {
}
{ } {
}
{ } {
}
}
{}{}{}
gilt.





}
{} {}

\inputaufgabe
{}
{





Zeige, dass die Aussage von

Aufgabe 5.21

nicht ohne die Voraussetzung gilt, dass die Abbildung polynomial ist.





}
{} {}





  \zwischenueberschrift{Aufgaben zum Abgeben}

\inputaufgabe
{3}
{





Wie viele
\definitionsverweis {Monome}{}{}
vom
\definitionsverweis {Grad}{}{}
$d$ gibt es im Polynomring in einer, in zwei und in drei Variablen?





}
{} {}

\inputaufgabe
{3}
{





Wende den Beweis zu

Satz 5.4

auf die algebraische Kurve an, die zur rationalen Funktion

\mavergleichskettedisp
{\vergleichskette
{Y
}
{ =} { { \frac{  X^2-2X }{ X^2-1 } }
}
{ } {
}
{ } {
}
{ } {
}
}
{}{}{}
gehört.





}
{} {}

\inputaufgabe
{3}
{





Betrachte die Abbildung
\maabbeledisp {} { {\mathbb A}^{2}_{ K } } { {\mathbb A}^{2}_{ K }
} {(x,y)} {(x,xy)
} {.}
Bestimme das Bild und die Fasern dieser Abbildung.





}
{} {}

\inputaufgabe
{3}
{





Betrachte das \stichwort { Ellipsoid} {}

\mavergleichskettedisp
{\vergleichskette
{ E
}
{ =} { V \left(  2x^2+3y^2+4z^2-5  \right)
}
{ =} { \left\{  (x,y,z)   \mid   2x^2+3y^2+4z^2 = 5         \right\}
}
{ } {
}
{ } {
}
}
{}{}{.}
Finde eine affin-lineare Variablentransformation
\zusatzklammer {über $\R$} {} {} derart, dass das Bild von $E$ unter der Abbildung die \stichwort {Standardkugel} {}
\mathl{V \left(  x^2+y^2+z^2-1  \right)}{} ist.





}
{} {}

\bild{ \begin{center}
\includegraphics[width=5.5cm]{ \bildeinlesungpng {Elipsoid_trojosy321} {png} }
\end{center}
\bildtext {Ein Ellipsoid: In der algebraischen Geometrie ist damit die Oberfläche gemeint.} }
\bildlizenz { Elipsoid trojosy321.png } {} {Pajs} {cz. Wikipedia} {gemeinfrei} {}





\inputaufgabe
{}
{





Es seien $V$ und $\tilde{V}$
\definitionsverweis {affin-algebraische Mengen}{}{}
in ${\mathbb A}^{2}_{ K }$ zu

\mavergleichskette
{\vergleichskette
{ K
}
{ = }{ \Z/( 2 )
}
{  }{
}
{    }{
}
{   }{
}
}
{}{}{.}
Zeige, dass diese beiden Mengen genau dann
\definitionsverweis {affin-linear äquivalent}{}{}
sind, wenn sie die gleiche Anzahl besitzen.





}
{Zeige ebenso, dass dies bei
\mathl{K=\Z/(p)}{} für
\mathl{p \geq 3}{} und auch für
\mathl{{ {\mathbb A}_{  \Z/(2)  }^{ n  } }}{} für
\mathl{n \geq 3}{} nicht gilt.} {}
